\documentclass[10pt,a4paper]{amsart}
\usepackage{amsmath,amssymb,amsthm,graphicx}
% Retained minimal preamble for the Stacks Project chapter varieties.tex.
% Verbatim excerpt of the chapter preamble: Xy-pic setup (native preamble lines 14-18), the native theorem
% declarations with the mathematical macros (native preamble lines 154-217), and the project own no-op
% definitions of the reference/slogan/history display environments (native wrapper.tex lines 1-3) that the
% counted statement uses for its slogan. Copied mechanically; nothing added, removed or redefined.
% For commutative diagrams we use Xy-pic
\usepackage[all]{xy}

% We use 2cell for 2-commutative diagrams.
\xyoption{2cell}
% Theorem environments.
%
\theoremstyle{plain}
\newtheorem{theorem}[subsection]{Theorem}
\newtheorem{proposition}[subsection]{Proposition}
\newtheorem{lemma}[subsection]{Lemma}

\theoremstyle{definition}
\newtheorem{definition}[subsection]{Definition}
\newtheorem{example}[subsection]{Example}
\newtheorem{exercise}[subsection]{Exercise}
\newtheorem{situation}[subsection]{Situation}

\theoremstyle{remark}
\newtheorem{remark}[subsection]{Remark}
\newtheorem{remarks}[subsection]{Remarks}

\numberwithin{equation}{subsection}

% Macros
%
\def\lim{\mathop{\mathrm{lim}}\nolimits}
\def\colim{\mathop{\mathrm{colim}}\nolimits}
\def\Spec{\mathop{\mathrm{Spec}}}
\def\Hom{\mathop{\mathrm{Hom}}\nolimits}
\def\Ext{\mathop{\mathrm{Ext}}\nolimits}
\def\SheafHom{\mathop{\mathcal{H}\!\mathit{om}}\nolimits}
\def\SheafExt{\mathop{\mathcal{E}\!\mathit{xt}}\nolimits}
\def\Sch{\mathit{Sch}}
\def\Mor{\mathop{\mathrm{Mor}}\nolimits}
\def\Ob{\mathop{\mathrm{Ob}}\nolimits}
\def\Sh{\mathop{\mathit{Sh}}\nolimits}
\def\NL{\mathop{N\!L}\nolimits}
\def\CH{\mathop{\mathrm{CH}}\nolimits}
\def\proetale{{pro\text{-}\acute{e}tale}}
\def\etale{{\acute{e}tale}}
\def\QCoh{\mathit{QCoh}}
\def\Ker{\mathop{\mathrm{Ker}}}
\def\Im{\mathop{\mathrm{Im}}}
\def\Coker{\mathop{\mathrm{Coker}}}
\def\Coim{\mathop{\mathrm{Coim}}}

% Boxtimes
%
\DeclareMathSymbol{\boxtimes}{\mathbin}{AMSa}{"02}

%
% Macros for moduli stacks/spaces
%
\def\QCohstack{\mathcal{QC}\!\mathit{oh}}
\def\Cohstack{\mathcal{C}\!\mathit{oh}}
\def\Spacesstack{\mathcal{S}\!\mathit{paces}}
\def\Quotfunctor{\mathrm{Quot}}
\def\Hilbfunctor{\mathrm{Hilb}}
\def\Curvesstack{\mathcal{C}\!\mathit{urves}}
\def\Polarizedstack{\mathcal{P}\!\mathit{olarized}}
\def\Complexesstack{\mathcal{C}\!\mathit{omplexes}}
% \Pic is the operator that assigns to X its picard group, usage \Pic(X)
% \Picardstack_{X/B} denotes the Picard stack of X over B
% \Picardfunctor_{X/B} denotes the Picard functor of X over B
\def\Pic{\mathop{\mathrm{Pic}}\nolimits}
\def\Picardstack{\mathcal{P}\!\mathit{ic}}
\def\Picardfunctor{\mathrm{Pic}}
\def\Deformationcategory{\mathcal{D}\!\mathit{ef}}
\newenvironment{reference}{}{}
\newenvironment{slogan}{}{}
\newenvironment{history}{}{}
\usepackage{cleveref}
\crefname{originallemma}{original lemma}{original lemmas}
\Crefname{originallemma}{Original lemma}{Original lemmas}
\makeatletter
\newlabel{lemma-chi-tensor-finite}{{0AYT}{source}{}{sourcepointer-1}{}}
\newlabel{lemma-chi-tensor-finite@cref}{{[originallemma][0][]0AYT}{[1][0][]source}{}{}{}}
\makeatother
\crefname{originallemma}{original lemma}{original lemmas}
\Crefname{originallemma}{Original lemma}{Original lemmas}
\makeatletter
\newlabel{lemma-intersection-number-in-terms-of-components}{{0BES}{source}{}{sourcepointer-2}{}}
\newlabel{lemma-intersection-number-in-terms-of-components@cref}{{[originallemma][0][]0BES}{[1][0][]source}{}{}{}}
\makeatother
\crefname{originallemma}{original lemma}{original lemmas}
\Crefname{originallemma}{Original lemma}{Original lemmas}
\makeatletter
\newlabel{lemma-intersection-number-additive}{{0BER}{source}{}{sourcepointer-3}{}}
\newlabel{lemma-intersection-number-additive@cref}{{[originallemma][0][]0BER}{[1][0][]source}{}{}{}}
\makeatother
\crefname{originalreference}{original reference}{original references}
\Crefname{originalreference}{Original reference}{Original references}
\makeatletter
\newlabel{properties-definition-ample}{{01PS}{source}{}{sourcepointer-4}{}}
\newlabel{properties-definition-ample@cref}{{[originalreference][0][]01PS}{[1][0][]source}{}{}{}}
\makeatother
\crefname{originalreference}{original reference}{original references}
\Crefname{originalreference}{Original reference}{Original references}
\makeatletter
\newlabel{morphisms-lemma-finite-proper}{{01WN}{source}{}{sourcepointer-5}{}}
\newlabel{morphisms-lemma-finite-proper@cref}{{[originalreference][0][]01WN}{[1][0][]source}{}{}{}}
\makeatother
\crefname{originalreference}{original reference}{original references}
\Crefname{originalreference}{Original reference}{Original references}
\makeatletter
\newlabel{divisors-section-effective-Cartier-divisors}{{01WQ}{source}{}{sourcepointer-6}{}}
\newlabel{divisors-section-effective-Cartier-divisors@cref}{{[originalreference][0][]01WQ}{[1][0][]source}{}{}{}}
\makeatother
\crefname{originalreference}{original reference}{original references}
\Crefname{originalreference}{Original reference}{Original references}
\makeatletter
\newlabel{divisors-section-effective-Cartier-invertible}{{0C4S}{source}{}{sourcepointer-7}{}}
\newlabel{divisors-section-effective-Cartier-invertible@cref}{{[originalreference][0][]0C4S}{[1][0][]source}{}{}{}}
\makeatother
\crefname{originalreference}{original reference}{original references}
\Crefname{originalreference}{Original reference}{Original references}
\makeatletter
\newlabel{divisors-section-Noetherian-effective-Cartier}{{0B3Q}{source}{}{sourcepointer-8}{}}
\newlabel{divisors-section-Noetherian-effective-Cartier@cref}{{[originalreference][0][]0B3Q}{[1][0][]source}{}{}{}}
\makeatother
\crefname{originallemma}{original lemma}{original lemmas}
\Crefname{originallemma}{Original lemma}{Original lemmas}
\makeatletter
\newlabel{lemma-dimension-locally-algebraic}{{0A21}{source}{}{sourcepointer-9}{}}
\newlabel{lemma-dimension-locally-algebraic@cref}{{[originallemma][0][]0A21}{[1][0][]source}{}{}{}}
\makeatother
\crefname{originallemma}{original lemma}{original lemmas}
\Crefname{originallemma}{Original lemma}{Original lemmas}
\makeatletter
\newlabel{lemma-numerical-intersection-effective-Cartier-divisor}{{0BEU}{source}{}{sourcepointer-10}{}}
\newlabel{lemma-numerical-intersection-effective-Cartier-divisor@cref}{{[originallemma][0][]0BEU}{[1][0][]source}{}{}{}}
\makeatother
\title{Intersection numbers of ample invertible sheaves on a proper scheme are positive}
\author{The Stacks Project Authors}
\date{stacks-project:varieties.tex:local-snapshot-20261011; GNU FDL 1.2 or later}
\begin{document}
\maketitle
\noindent\textit{Editorial note.} All delivered mathematical text is the author's own, copied verbatim from the pinned Stacks Project chapter file varieties.tex; no statement or proof step was rewritten, reordered or removed.\par\medskip
\noindent\textit{Editorial note.} The counted span cites ten cross-chapter labels (including a definition and three section labels), each recorded in assembly.external\_references with its canonical tag as a navigation pointer.\par\medskip
\noindent\textit{Editorial note.} The proof cites no bibliography item inside the counted span.\par\medskip
\noindent\textit{Allowed standard prerequisites (editorial).}\begin{itemize}
\item Characterization of ample invertible sheaves by affine nonvanishing loci of sections. Stacks Project, Properties, Definition properties-definition-ample (tag 01PS). Produces the section s with X\_s affine after replacing L\_1 by a positive tensor power.
\item Effective Cartier divisors and the intersection number reduction. Stacks Project, Varieties, Lemmas lemma-numerical-intersection-effective-Cartier-divisor (tag 0BEU) and lemma-intersection-number-in-terms-of-components (tag 0BES). Rewrite the intersection number on the divisor Z(s) and reduce to the integral case.
\item Additivity of intersection numbers and finiteness of cohomology in dimension zero. Stacks Project, Varieties, Lemma lemma-intersection-number-additive (tag 0BER) and Lemma lemma-chi-tensor-finite (tag 0AYT). Allow replacing L\_1 by a positive tensor power and supply the base case.
\end{itemize}\medskip
\ifdefined\hypertarget\hypertarget{sourcepointer-1}{}\fi
\noindent\textit{Source reference.} Original lemma 0AYT: sources/stacks-chow/tags.txt, source lines 13801--13801. Cross-chapter standard prerequisite cited by the author inside the counted span; it is used as a pointer and its own proof is not part of this item.\par\medskip
\ifdefined\hypertarget\hypertarget{sourcepointer-2}{}\fi
\noindent\textit{Source reference.} Original lemma 0BES: sources/stacks-chow/tags.txt, source lines 14360--14360. Cross-chapter standard prerequisite cited by the author inside the counted span; it is used as a pointer and its own proof is not part of this item.\par\medskip
\ifdefined\hypertarget\hypertarget{sourcepointer-3}{}\fi
\noindent\textit{Source reference.} Original lemma 0BER: sources/stacks-chow/tags.txt, source lines 14359--14359. Cross-chapter standard prerequisite cited by the author inside the counted span; it is used as a pointer and its own proof is not part of this item.\par\medskip
\ifdefined\hypertarget\hypertarget{sourcepointer-4}{}\fi
\noindent\textit{Source reference.} Original reference 01PS: sources/stacks-chow/tags.txt, source lines 2241--2241. Cross-chapter standard prerequisite cited by the author inside the counted span; it is used as a pointer and its own proof is not part of this item.\par\medskip
\ifdefined\hypertarget\hypertarget{sourcepointer-5}{}\fi
\noindent\textit{Source reference.} Original reference 01WN: sources/stacks-chow/tags.txt, source lines 2488--2488. Cross-chapter standard prerequisite cited by the author inside the counted span; it is used as a pointer and its own proof is not part of this item.\par\medskip
\ifdefined\hypertarget\hypertarget{sourcepointer-6}{}\fi
\noindent\textit{Source reference.} Original reference 01WQ: sources/stacks-chow/tags.txt, source lines 2491--2491. Cross-chapter standard prerequisite cited by the author inside the counted span; it is used as a pointer and its own proof is not part of this item.\par\medskip
\ifdefined\hypertarget\hypertarget{sourcepointer-7}{}\fi
\noindent\textit{Source reference.} Original reference 0C4S: sources/stacks-chow/tags.txt, source lines 15237--15237. Cross-chapter standard prerequisite cited by the author inside the counted span; it is used as a pointer and its own proof is not part of this item.\par\medskip
\ifdefined\hypertarget\hypertarget{sourcepointer-8}{}\fi
\noindent\textit{Source reference.} Original reference 0B3Q: sources/stacks-chow/tags.txt, source lines 13973--13973. Cross-chapter standard prerequisite cited by the author inside the counted span; it is used as a pointer and its own proof is not part of this item.\par\medskip
\ifdefined\hypertarget\hypertarget{sourcepointer-9}{}\fi
\noindent\textit{Source reference.} Original lemma 0A21: sources/stacks-chow/tags.txt, source lines 12689--12689. Cross-chapter standard prerequisite cited by the author inside the counted span; it is used as a pointer and its own proof is not part of this item.\par\medskip
\ifdefined\hypertarget\hypertarget{sourcepointer-10}{}\fi
\noindent\textit{Source reference.} Original lemma 0BEU: sources/stacks-chow/tags.txt, source lines 14362--14362. Cross-chapter standard prerequisite cited by the author inside the counted span; it is used as a pointer and its own proof is not part of this item.\par\medskip
\begin{lemma}
\label{lemma-ample-positive}
Let $k$ be a field. Let $X$ be proper over $k$. Let $Z \subset X$ be
a closed subscheme of dimension $d$. If $\mathcal{L}_1, \ldots, \mathcal{L}_d$
are ample, then $(\mathcal{L}_1 \cdots \mathcal{L}_d \cdot Z)$ is positive.
\end{lemma}

\begin{proof}
We will prove this by induction on $d$. The case $d = 0$
follows from Lemma \ref{lemma-chi-tensor-finite}. Assume $d > 0$.
By Lemma \ref{lemma-intersection-number-in-terms-of-components}
we may assume that $Z$ is an integral closed subscheme.
In fact, we may replace $X$ by $Z$ and $\mathcal{L}_i$
by $\mathcal{L}_i|_Z$ to reduce to the case $Z = X$ is a
proper variety of dimension $d$.
By Lemma \ref{lemma-intersection-number-additive}
we may replace $\mathcal{L}_1$ by a positive tensor power.
Thus we may assume there exists a nonzero section
$s \in \Gamma(X, \mathcal{L}_1)$
such that $X_s$ is affine (here we use the definition of
ample invertible sheaf, see
Properties, Definition \ref{properties-definition-ample}).
Observe that $X$ is not affine because proper and affine
implies finite (Morphisms, Lemma \ref{morphisms-lemma-finite-proper})
which contradicts $d > 0$. It follows that $s$ has a nonempty vanishing
scheme $Z(s) \subset X$. Since $X$ is a variety, $s$ is a regular section
of $\mathcal{L}_1$, so $Z(s)$ is an effective Cartier divisor,
thus $Z(s)$ has codimension $1$ in $X$, and
hence $Z(s)$ has dimension $d - 1$ (here we use material from
Divisors, Sections \ref{divisors-section-effective-Cartier-divisors},
\ref{divisors-section-effective-Cartier-invertible}, and
\ref{divisors-section-Noetherian-effective-Cartier}
and from dimension theory as in Lemma \ref{lemma-dimension-locally-algebraic}).
By Lemma \ref{lemma-numerical-intersection-effective-Cartier-divisor}
we have
$$
(\mathcal{L}_1 \cdots \mathcal{L}_d \cdot X) =
(\mathcal{L}_2 \cdots \mathcal{L}_d \cdot Z(s))
$$
By induction the right hand side is positive and the proof is complete.
\end{proof}

\end{document}
